In the last decade, computing systems have become an integral part of our
society, with all kinds of devices being physically present in the world to
perform a multitude of tasks, ranging from everyday activities like browsing
social media and playing video games to sensitive operations such as storing
bank and health data, controlling robots in industry systems, or driving
vehicles. Furthermore, boosted by the advent of cloud computing and the
\ac{IoT}, many applications that originally were composed of a single
monolithical unit have been split into several parts that are geographically
distributed and interconnected by a network, continuously exchanging data to
provide their services. Making traditional systems and appliances intelligent
does have several advantages that can, among others, improve quality of life and
reduce costs and waste. In agriculture, for example, devices such as smart
irrigation systems can dramatically improve efficiency by saving large amounts
of water, which not only reduces costs but also has a positive impact on the
environment~\cite{elsayed2021smartFarming}. Similarly, drones can be used to
monitor the growth of crops and create prescription maps to optimize the
application of products such as fertilizers or pesticides~\cite{inoue2020drone}.
From another point of view, assisted and automated driving can increase the
safety on the road and significantly reduce the number of accidents and
fatalities~\cite{nhtsaVehicleSafety}.

However, with an increased number of connected devices and applications, several
concerns have arisen. On the one hand, exposing services and appliances on the
Internet has made it much easier for attackers to breach into them, which has
led to several incidents in the past: For example, in 2016 the Mirai botnet
caused the largest distributed denial of service attack ever recorded, thanks to
a malware that spread over thousands of vulnerable \ac{IoT} devices connected to
the Internet~\cite{miraiBotnet}. More recently, in 2021, the Colonial Pipeline
oil company was victim of a ransomware attack that caused the disruption of the
fuel supply chain in the United States for several days, due to an insecure VPN
interface~\cite{colonialPipelineAttack}. On the other hand, applications might
process sensitive information such as health data or bank details, which raises
some interesting privacy challenges: Guidelines like the NIST cybersecurity
framework~\cite{nistcsf} and regulations such as the \ac{GDPR} in Europe and the
\ac{CCPA} in the United States have been adopted to ensure that companies comply
with data protection principles to prevent leaks of users' personal information.
Nevertheless, the number of data breaches in the last years still shows a rising
trend\footnote{\url{https://www.statista.com/statistics/273550/data-breaches-recorded-in-the-united-states-by-number-of-breaches-and-records-exposed/}}.
To deal with this increased attack surface, several defense mechanisms have been
proposed over the years, some of which have been adopted as standard
countermeasure in all systems. Cryptographic techniques such as symmetric and
asymmetric encryption are now the standard way to protect the confidentiality
and integrity of data, both at rest (e.g., when stored on a hard drive) and in
transit (e.g., during communication between two services). Besides, a mechanism
to manage identities and certificates, such as a \ac{PKI}~\cite{pki}, can
provide additional authenticity guarantees, e.g., to make sure that a user is
communicating with a legitimate server. Instead, protecting data ``in use'',
i.e., while being processed within a device, is still an open research area that
has gained considerable interest in recent years, as workloads are being
offloaded more and more from local, trusted machines to remote servers managed
by third parties and whose resources are shared among multiple processes.

In this Ph.D. thesis, we focus on a specific solution called
\acp{TEE}~\cite{schneider2022hardware}, a hardware-based mechanism that allows
the processor to isolate portions of code and data during execution on a system,
preventing even privileged code from accessing them. Despite being a relatively
new technology, \acp{TEE} are getting a lot of traction in recent years,
especially since the release of Intel \ac{SGX} in 2015~\cite{sgx} that made this
technology available in commodity x86 processors. Nowadays, \acp{TEE} are more
popular in the context of cloud computing where physical resources are shared
among multiple tenants and strong workload isolation is needed. Besides,
\acp{TEE} can shield data even from the infrastructure provider, which allows
moving applications to the cloud while still retaining the privacy of user data,
which is crucial for highly sensitive tasks. This paradigm is known as
\emph{confidential computing}, and it is the main focus of recent architectures
like AMD \sevsnp{}~\cite{AmdSevSnpWhitepaper} and Intel
\ac{TDX}~\cite{intelTdxWhitepaper}, which became available in 2021 and 2023,
respectively. \acp{TEE} can provide additional security guarantees to existing
applications and enable novel use cases and protocols that were not possible
before~\cite{caseConfidentialComputing}.

However, integrating \acp{TEE} into distributed applications poses several
technical and operational challenges. Indeed, running an application inside a
\ac{TEE} requires careful redesign and some code changes to isolate the
sensitive code and data that need to be protected. Even though recent \acp{TEE}
such as AMD \sevsnp{} and Intel \ac{TDX} can run unmodified binaries, developers
might still need to make some adjustments to avoid using resources that may
still be influenced by attackers, either directly or indirectly. Besides, there
must be a mechanism in place to actually \emph{verify} that the application is
running in such an isolated environment without being compromised, so that users
can \emph{trust} that any data exchanged with the application will be not
accessible by unauthorized parties. This process is called \emph{remote
attestation}
and, while it is a feature offered by all \acp{TEE}, it requires some
development effort and expertise to implement it correctly. Remote attestation
can also be mediated by third party services or even cloud providers, though
this requires an in-depth analysis of the threat model and the trust
relationships between different entities.

In this dissertation, we advance the state of the art in \ac{TEE} research on
two different dimensions: First, we demonstrate how integrating trusted
computing technologies into common protocols and applications can provide
several advantages in terms of security, and in some cases even increased
efficiency and scalability compared to traditional appraches. To do so, we
showcase two novel use cases in the context of distributed \ac{IoT} and \ac{V2X}
settings. Second, we uncover some operational challenges and limitations that
should be taken into account when using \acp{TEE} in practice. In particular, we
discuss how getting a stable notion of time inside a \ac{TEE} is a non-trivial
task, and we shed some light on trust relationships between tenants and cloud
providers when deploying confidential computing solutions to public clouds.